\section*{Preface}
\addcontentsline{toc}{section}{Preface}

I did not seek this work. It found me, as such things always do, in the space between the tide and the truth.

The waters of the Amaranthine do not give up their dead easily. I learned this as a young lieutenant, pulling bodies from the wreck of the \emph{Mistral}---a Kahfagian cutter that had run afoul of something that had no business being in the shallows. The sailors had not drowned. I will not describe what I saw, except to say that I spent the next three years learning that the world is older and hungrier than the Admiralty's charts admit.

The attack on the cultists of Khemesh---that ill-fated expedition that claimed so many good men and women---left me broken on a beach on the Dhaharan coast. I washed ashore with salt in my lungs and a silver mirror in my hand that had not been there when I went under. When I opened my eyes, I saw not the sky, but a reflection of something that had been watching me for a very long time.

The Witness does not speak. It shows.

What follows is not a grimoire in the traditional sense. I have no interest in teaching the eager how to summon what should remain sealed. This is a casebook---a compendium of observations, autopsies, interrogations, and the hard-won wisdom of twenty years walking the edge between the living and the hungry.

I have organized this work by category, though I caution the reader that the categories are not clean. A ghost may become a ghoul, given enough hunger and the right soil. A demon may wear a vampire's mask. The lycanthrope's curse is kin to the plague-bearer's affliction, and both answer to Malachai's silver tongue.

The Witness demands truth. I have given it as I have found it.

\begin{flushright}
Saikou Ira\\
Port Shed, Dhaharan Coast\\
The ink is ash and silver. The mirror is clean.
\end{flushright}

\subsection*{A Note on Patrons and Tools}

I am an Invoker. I carry no Codex carved in stone, no Thiasos bound by oath. I carry Symbols---physical anchors to Patrons I have learned to bargain with, each for a specific purpose. Over years of walking the edge, I have found four powers that answer when I call.

\begin{description}
\item[The Witness] My first and most constant patron. It shows me what is hidden, reveals the truth beneath the lie, and guides my investigations. When I place a silver mirror, I invoke The Witness. Its domain is revelation, memory, and the unbearable clarity of the seen.

\item[The Traveler] Lord of roads, waystations, and the space between departure and arrival. I call on The Traveler when I must escape a trap, find a path that does not appear on any map, or simply arrive before my pursuers. Its gifts are subtle---a shortcut that should not exist, a door that opens to the right alley, a helpful stranger who speaks the language of the road.

\item[Raéyn] Mistress of the sea, storms, and the tide that does not forgive. I invoke Raéyn when I am desperate---when the water rises, when the ship is sinking, when the creature from the deep has me in its grip. Her favor is fickle, but her wrath is absolute.

\item[The Sentinel] (Some call it the Sealed Gate.) Patron of thresholds, wards, and the boundary between what belongs and what must be kept out. I call on The Sentinel to exorcise spirits, to seal a room against possession, to draw the line that nothing may cross. Its power is that of the \emph{witnessed barrier}---a line of salt blessed by a third party, a mirror facing a door, a vow spoken in the presence of two watchers.
\end{description}

In the chapters that follow, I will describe---when relevant---how I have used these patrons in particular cases. The rituals are not provided as instructions. They are provided as \emph{evidence} that such bargains are possible. Whether you choose to walk the same roads is your own affair.

\vspace{0.5cm}
\noindent\emph{Saikou Ira}

\section*{Preface}
\addcontentsline{toc}{section}{Preface}

I did not seek this work. It found me, as such things always do, in the space between the tide and the truth.

The waters of the Amaranthine do not give up their dead easily. I learned this as a young lieutenant, pulling bodies from the wreck of the \emph{Mistral}---a Kahfagian cutter that had run afoul of something that had no business being in the shallows. The sailors had not drowned. I will not describe what I saw, except to say that I spent the next three years learning that the world is older and hungrier than the Admiralty's charts admit.

The attack on the cultists of Khemesh---that ill-fated expedition that claimed so many good men and women---left me broken on a beach on the Dhaharan coast. I washed ashore with salt in my lungs and a silver mirror in my hand that had not been there when I went under. When I opened my eyes, I saw not the sky, but a reflection of something that had been watching me for a very long time.

The Witness does not speak. It shows.

What follows is not a grimoire in the traditional sense. I have no interest in teaching the eager how to summon what should remain sealed. This is a casebook---a compendium of observations, autopsies, interrogations, and the hard-won wisdom of twenty years walking the edge between the living and the hungry.

I have organized this work by category, though I caution the reader that the categories are not clean. A ghost may become a ghoul, given enough hunger and the right soil. A demon may wear a vampire's mask. The lycanthrope's curse is kin to the plague-bearer's affliction, and both answer to Malachai's silver tongue.

The Witness demands truth. I have given it as I have found it.

\begin{flushright}
Saikou Ira\\
Port Shed, Dhaharan Coast\\
The ink is ash and silver. The mirror is clean.
\end{flushright}

\subsection*{A Note on Patrons and Tools}

I am an Invoker. I carry no Codex carved in stone, no Thiasos bound by oath. I carry Symbols---physical anchors to Patrons I have learned to bargain with, each for a specific purpose. Over years of walking the edge, I have found four powers that answer when I call.

\begin{description}
\item[The Witness] My first and most constant patron. It shows me what is hidden, reveals the truth beneath the lie, and guides my investigations. When I place a silver mirror, I invoke The Witness. Its domain is revelation, memory, and the unbearable clarity of the seen.

\item[The Traveler] Lord of roads, waystations, and the space between departure and arrival. I call on The Traveler when I must escape a trap, find a path that does not appear on any map, or simply arrive before my pursuers. Its gifts are subtle---a shortcut that should not exist, a door that opens to the right alley, a helpful stranger who speaks the language of the road.

\item[Raéyn] Mistress of the sea, storms, and the tide that does not forgive. I invoke Raéyn when I am desperate---when the water rises, when the ship is sinking, when the creature from the deep has me in its grip. Her favor is fickle, but her wrath is absolute.

\item[The Sentinel] (Some call it the Sealed Gate.) Patron of thresholds, wards, and the boundary between what belongs and what must be kept out. I call on The Sentinel to exorcise spirits, to seal a room against possession, to draw the line that nothing may cross. Its power is that of the \emph{witnessed barrier}---a line of salt blessed by a third party, a mirror facing a door, a vow spoken in the presence of two watchers.
\end{description}

In the chapters that follow, I will describe---when relevant---how I have used these patrons in particular cases. The rituals are not provided as instructions. They are provided as \emph{evidence} that such bargains are possible. Whether you choose to walk the same roads is your own affair.

\vspace{0.5cm}
\noindent\emph{Saikou Ira}

